\addchap{Notation}

One notation is used throughout the book. Where a chapter needs a local symbol that clashes with this list, the chapter says so; where a source paper uses a different notation, the chapter's notes translate it.

\newcommand{\notationrow}[2]{#1 & #2\\}
\newcommand{\notationhead}[1]{\multicolumn{2}{@{}l}{\textsc{\color{rubric}#1}}\\[0.3ex]}
\noindent{\small
\renewcommand{\arraystretch}{1.1}%
\begin{tabular}{@{}>{\raggedright\arraybackslash}p{32mm}>{\raggedright\arraybackslash}p{83.5mm}@{}}
\notationhead{runs and comparisons}
\notationrow{run}{one training of one network, from initialization to the end; the experimental unit}
\notationrow{$Y$}{the outcome of a run: a test error, an accuracy, or a loss}
\notationrow{$n$}{the number of runs (replicates, seeds) in one arm of an experiment}
\notationrow{arms $1$, $2$}{the two arms of a comparison (in words: method A and method B, or treatment and control)}
\notationrow{$\mu_1$, $\mu_2$, $\Delta$}{the mean outcomes of the two arms, and the effect $\Delta=\mu_1-\mu_2$}
\notationrow{$\sigma$; $\sigma_1$, $\sigma_2$}{the standard deviation of the run distribution; of each arm, when the two differ}
\notationrow{trial}{one hyperparameter configuration, trained once or several times}
\notationrow{$h$, $\cH$, $k$}{a hyperparameter configuration, the search space, and a tuning budget of $k$ trials}
\addlinespace
\notationhead{statistics}
\notationrow{$\bar Y$, $S$}{the sample mean and the sample standard deviation (divisor $n-1$)}
\notationrow{$\SE$}{the standard error; $S/\sqrt n$ for a mean}
\notationrow{$\Prob$, $\E$, $\Var$, $\Cov$, $\Corr$}{probability, expectation, variance, covariance, and correlation}
\notationrow{$\Phi$, $z_q$}{the standard normal distribution function and its $q$-quantile}
\notationrow{$t_{\nu,q}$}{the $q$-quantile of Student's $t$ distribution with $\nu$ degrees of freedom}
\notationrow{$\alpha$}{a significance level; confidence intervals are at level $1-\alpha$, usually $95$ percent}
\notationrow{power}{the probability that a test detects a given effect, written in words (power $0.8$)}
\notationrow{\pmval{71.3}{0.4}}{a mean and the half-width of its $95$ percent interval}
\notationrow{$Y_{(1)}\le\dots\le Y_{(n)}$}{order statistics}
\notationrow{$M_k$, $e_k$}{the best of $k$ draws, and the expected maximum of $k$ independent standard normals}
\notationrow{$\tau$, $\IQM$, $\med$}{Kendall's rank correlation, the interquartile mean, and the median}
\notationrow{$\cN$, $\Bin$, $\iid$}{normal and binomial laws; $Y_1,\dots,Y_n\iid F$ means that the $Y_i$ are independent, each with law $F$}
\end{tabular}

\medskip\noindent
\begin{tabular}{@{}>{\raggedright\arraybackslash}p{32mm}>{\raggedright\arraybackslash}p{83.5mm}@{}}
\notationhead{training}
\notationrow{$\theta$}{the parameters of a network}
\notationrow{$\Loss(\theta)$, $\Lhat(\theta)$}{population loss and training loss}
\notationrow{$\eta$, $B$, $T$}{learning rate, batch size, and number of optimizer steps}
\notationrow{$\lambda$}{weight decay; eigenvalues always carry a subscript, as in $\lambda_{\max}$}
\notationrow{$\beta_1$, $\beta_2$, $\epsilon$}{the momentum, second-moment, and stability parameters of Adam}
\notationrow{$\nabla^2\Lhat(\theta)$}{the Hessian of the training loss; $\lambda_{\max}$ its largest eigenvalue}
\addlinespace
\notationhead{scale}
\notationrow{$N$, $D$, $C$}{parameters, training data (tokens or examples seen), and compute in floating-point operations; $C\approx 6ND$ for dense networks}
\notationrow{$\Loss(N,D)$}{the parametric law $E+A_NN^{-\alpha_N}+A_DD^{-\alpha_D}$, with irreducible loss $E$}
\notationrow{$a$, $b$}{the exponents of compute-optimal allocation: $N_{\mathrm{opt}}\propto C^{a}$, $D_{\mathrm{opt}}\propto C^{b}$}
\addlinespace
\notationhead{evaluation}
\notationrow{$m$}{the number of test items}
\notationrow{$\hat p$}{an estimated accuracy, with standard error $\sqrt{\hat p(1-\hat p)/m}$}
\notationrow{$J$, $j$}{the number of tasks in a suite, and a task}
\addlinespace
\notationhead{status of a statement}
\notationrow{$\blacksquare$}{end of a complete proof}
\notationrow{proof sketch}{an argument that omits steps, with a pointer to where they are}
\notationrow{$\Diamond$}{end of a heuristic derivation: an argument that is not a proof}
\notationrow{finding}{an empirical result as reported, with its source and its size}
\end{tabular}\par}
